\chapter{91}\label{section-91}

\section{Mailand (IT), Montag, 10.
September}\label{mailand-it-montag-10.-september}

Commissario Alessandro Mantovani

Emma steht im Türrahmen meines Büros. Sie hebt fragend die Hand.

«Ja, bitte, Emma,» ich weise auf einen Platz mir gegenüber am
Schreibtisch, «alles klar für dich?»

«Ja,» und während sie nickt, sehe ich in ihrem Gesicht diese
Konzentration, die ich an ihr mag.

«Gut, nochmal kurz, ich führe die Befragung durch, du gibst ihr
Sicherheit und bringst dich bitte ein, wenn es dir wichtig erscheint,
wie immer, mit vorgängigem Blickkontakt?»

Emma schliesst kurz die Augen. Sie hat mich verstanden. Dann weist sie
zur Tür.

«Ich hole sie.»

«Buongiorno, Signora Ricci.»

Ich reiche der jungen Frau meine Hand. Ihre ist feucht. Der kräftige
Druck erstaunt mich. Emma rückt ihr den Stuhl zurecht. Die Augen der
Frau flattern. Sie setzt sich ganz vorne an die Stuhlkante.

«Danke, dass Sie sich gemeldet haben, Signora Ricci.»

Ich suche ihren Blick. Wähle eine beruhigende Stimmlage.

«Zu ihren Personalien: Valentina Ricci, wohnhaft in Mailand.»

Sie hält meinem Blick stand und nickt die Angaben mit kleinen
Kopfbewegungen ab.

Jetzt spreche ich langsamer und einen Deut tiefer.

«Sie sind bei uns an einem sicheren Ort, Signora. Sie können frei
sprechen. Bitte erzählen Sie uns noch einmal der Reihe nach, was Sie
bereits am Telefon sagten, okay?»

Die Frau legt nun ihre Hände ineinander. Sie rutscht tatsächlich auf dem
Stuhl ein Stück nach hinten und atmet tief ein. Dennoch sehe ich die
Angst in ihren Augen.

Ich hoffe, sie wagt es, sich einzulassen.

Emma dreht sich aufmunternd zu ihr hin. Ich lehne mich zurück,
überschlage die Beine und blicke bewusst auf meine Notizen auf dem
Schreibtisch.

«Ich, ehm, ich musste sie anschwindeln,» beginnt sie, «es tut mir leid.»

Ich nicke. Mit einem zusätzlichen Lächeln will ich ihr Verständnis
signalisieren. Emma verhält sich synchron zu mir.

Je länger die Frau erzählt, desto weiter rutscht sie auf dem Stuhl nach
hinten. Schliesslich berührt ihr Rücken die Stuhllehne.

«Ja, ich habe Kontakt mit Louis Vincent, er, ehm, er hat mich
eindringlich gebeten, niemandem davon zu erzählen, er ist mir ein guter
Freund, darum, aber,» ihre Stimme wird erstaunlich fest, «aber, was am
\emph{Roccia} ganz genau geschehen ist, weiss ich nicht. Das müssen Sie
mir glauben.»

Sie blickt zwischen Emma und mir hin und her. Auf ihrem Gesicht erkenne
ich diesen flehenden Ausdruck, der mir bekannt ist. In der Regel drückt
er Aufrichtigkeit aus. Ich kneife die Augen zusammen und betrachte sie.
Eigentlich müsste jemand die junge Frau jetzt in den Arm nehmen.
Vielleicht liegt es am Alter, dass ich zunehmend rührseliger werde? Ich
wische den Gedanken weg und schlüpfe zurück in meine Rolle.

Emma berührt die Frau am Arm. Dann schaut sie mich kurz an.

Während Emma spricht, werden die Augen der jungen Frau nach und nach
glasiger.

«Signora Ricci, bitte beschreiben Sie den Übergriff der beiden Männer so
exakt, wie es Ihnen möglich ist. Jedes Detail kann uns weiterhelfen. Wir
verstehen Ihre Angst sehr gut. Die Männer haben Sie massiv bedroht.»
